\BourbakiBackMatter{参考文献}

{[}1{]} E. ARTIN -- ``Zur Theorie der hypercomplexen Zahlen'', Abh. math. Sem. Univ. Hamburg 5 (1928), p.~251--260; Collected papers, Reading (Addison Wesley), 1965, p.~307--316.

{[}2{]} R. BRAUER -- ``Über Systeme hypercomplexer Zahlen'', Math. Zeitschr. 30 (1929), p.~79--107; Collected papers, vol.~I, Cambridge, Massachusetts (The MIT Press), 1980, p.~40--68.

{[}3{]} W. BURNSIDE - ``论任意线性替换群的可约性条件'', Proc. Lond. Math. Soc. 3 (1905), p.~430-434.

{[}4{]} É. CARTAN -- ``Les groupes bilinéaires et les systèmes de nombres complexes'', Ann. Fac. Sc. Toulouse 12 (1898), p.~1--99; Œuvres complètes, partie II, vol.~I, Paris (Gauthier-Villars), 1952, p.~7--105.

{[}5{]} \_\_\_\_ ``Nombres complexes'', Encycl. Sci. Math., éd. française 15 (1908); Œuvres complètes, partie II, vol.~I, Paris (Gauthier-Villars), 1952, p.~107--246.

{[}6{]} A. CAYLEY -- ``论代数偶'', Phil. Mag. (3) 27 (1845), p.~38--40; Collected mathematical papers, t. I, Cambridge, 1889, p.~128--131.

{[}7{]} \_\_\_\_ ``论依赖于方程 \(\theta^{n}=1\) 的群论'', Phil. Mag. (4) 7 (1854), p.~40--47; Collected mathematical papers, t. II, Cambridge, 1889, p.~123--130.

{[}8{]} \_\_\_\_``矩阵理论专论'', Phil. Trans. R. Soc. Lond. 148 (1858), p.~17--37; Collected mathematical papers, t. II, Cambridge, 1889, p.~475--496.

{[}9{]} W. CLIFFORD -- ``双四元数初步概要'', Proc. Lond. Math. Soc. 4 (1873), p.~381--395; Mathematical papers, London (Macmillan), 1882, p.~181--200.

{[}10{]} \_\_\_\_ ``论几何代数的分类'', (1877), Mathematical papers, London (Macmillan), 1882, p.~397--401.

{[}11{]} \_\_\_\_``格拉斯曼扩张代数的应用'', Amer. Journ. of Math. 1 (1878), p.~350--358; Mathematical papers, London (Macmillan), 1882, p.~266--276.

{[}12{]} R. DEDEKIND -- ``Zur Theorie der aus n Haupteinheiten gebildeten komplexen Grössen'', Gött. Nachr. (1885), p.~141--159; Gesammelte mathematische Werke, t. II, Braunschweig (Vieweg), 1931--32, p.~1--19.

{[}13{]} \_\_\_\_ ``Aus Briefen an Frobenius'', (1886), Gesammelte mathematische Werke, t. II, Braunschweig (Vieweg), 1931--32, p.~414--442.

{[}14{]} \_\_\_\_ ``Über eine Erweiterung des Symbols (a, b) in der Theorie der Moduln'', Gött. Nachr. (1895), p.~183--188; Gesammelte mathematische Werke, t. II, Braunschweig (Vieweg), 1931--32, p.~59--86.

{[}15{]} \_\_\_\_``Über Gruppen, deren sämtliche Teiler Normalteiler sind'', Math. Ann. 48 (1897), p.~548--561; Gesammelte mathematische Werke, t. II, Braunschweig (Vieweg), 1931--32, p.~87--102.

{[}16{]} M. DEURING - Algebren, Erg. der Math., vol.~4, Berlin (Springer), 1935.

{[}17{]} L. DICKSON - ``线性结合代数与阿贝尔方程'', Trans. Amer. Math. Soc. 15 (1914), p.~31-46; Mathematical papers, vol.~II, New York (Chelsea), 1975, p.~367-382.

{[}18{]} G. FROBENIUS -- ``Über lineare Substitutionen und bilineare Formen'', Crelle's J. 84 (1878), p.~1--63; Gesammelte Abhandlungen, Band I, Berlin (Springer), 1968, p.~343--405.

{[}19{]} \_\_\_\_ ``Über die Primfactoren der Gruppendeterminante'', Berliner Sitzungsber. (1896), p.~1343--1382; Gesammelte Abhandlungen, Band III, Berlin (Springer), 1968, p.~38--77.

{[}20{]} \_\_\_\_ ``Über Gruppencharaktere'', Berliner Sitzungsber. (1896), p.~985--1021; Gesammelte Abhandlungen, Band III, Berlin (Springer), 1968, p.~1--37.

{[}21{]} \_\_\_\_``Über die Darstellung der endlichen Gruppen durch lineare Substitutionen'', Berliner Sitzungsber. (1897), p.~944--1015; Gesammelte Abhandlungen, Band III, Berlin (Springer), 1968, p.~82--103.

{[}22{]} \_\_\_\_ ``Theorie der hypercomplexen Grössen'', Berliner Sitzungsber. (1903), p.~504--537 与 634--645; Gesammelte Abhandlungen, Band III, Berlin (Springer), 1968, p.~284--329.

{[}23{]} H. GRASSMANN -- ``Sur les différents genres de multiplication'', Crelle's J. 49 (1855), p.~123--141; Gesammelte math. und phys. Werke, t. II, Leipzig (Teubner), 1904, p.~199--217.

{[}24{]} \_\_\_\_ ``Die Ausdehnungslehre von 1862'', (1862), Gesammelte math. und phys. Werke, t. I, Leipzig (Teubner), 1896, p.~1--383.

{[}25{]} W. HAMILTON - 四元数讲义, Dublin (Hodges and Smith), 1853.

{[}26{]} O. HÖLDER -- ``Zurückführung einer beliebigen algebraischen Gleichung auf eine Kette von Gleichungen'', Math. Ann. 34 (1889), p.~26--56.

{[}27{]} C. HOPKINS - ``具有左理想极小条件的环'', Ann. of Math. 40 (1939), p.~712-730.

{[}28{]} W. KRULL -- ``Über verallgemeinerte endliche Abelsche Gruppen'', Math. Zeitschr. 23 (1925), p.~161--196; Gesammelte Abhandlungen, vol.~1, Berlin (de Gruyter), 1999, p.~263--298.

{[}29{]} \_\_\_\_ ``Theorie und Anwendung der verallgemeinerten Abelschen Gruppen'', Sitzungsber. Heidelberger Akad. (1926), p.~1--32; Gesammelte Abhandlungen, vol.~1, Berlin (de Gruyter), 1999, p.~299--328.

{[}30{]} E. LAGUERRE -- ``Sur le calcul des systèmes linéaires'', J. de l'école polytechnique 42 (1867), p.~215--264; Œuvres, t. I, Paris (Gauthier-Villars).

{[}31{]} T. MOLIEN - ``Ueber Systeme höherer complexer Zahlen'', Math. Ann. 41 (1893), p.~83-156.

{[}32{]} \_\_\_\_``Ueber die Invarianten der linearen Substitutionsgruppen'', Berliner Sitzungsber. (1897), p.~1152--1156.

{[}33{]} E. NOETHER -- ``Abstrakter Aufbau der Idealtheorie in algebraischen Zahl- und Funktionenkörpern'', Math. Ann. 96 (1927), p.~26--61; Gesammelte Abhandlungen, Berlin (Springer), 1983, p.~493--528.

{[}34{]} \_\_\_\_ ``Hyperkomplexe Grössen und Darstellungstheorie'', Math. Zeitschr. 30 (1929), p.~641--692; Gesammelte Abhandlungen, Berlin (Springer), 1983, p.~563--614.

{[}35{]} \_\_\_\_ ``Nichtkommutative Algebren'', Math. Zeitschr. 37 (1933), p.~514--541; Gesammelte Abhandlungen, Berlin (Springer), 1983, p.~642--669.

{[}36{]} E. NOETHER et R. BRAUER -- ``Über minimale Zerfällungskörper irreduzibler Darstellungen'', Berliner Sitzungsber. (1927), p.~221--228; Emmy Noether Gesammelte Abhandlungen, Berlin (Springer), 1983, p.~552--559.

{[}37{]} E. NOETHER et W. SCHMEIDLER -- ``Moduln in nichtkommutativen Bereichen'', Math. Zeitschr. 8 (1920), p.~1--35; Emmy Noether Gesammelte Abhandlungen, Berlin (Springer), 1983, p.~318--352.

{[}38{]} B. PEIRCE - ``线性结合代数'', Amer. Journ. of Math. 4 (1881), p.~97-221.

{[}39{]} C. PEIRCE - ``论除法无歧义的代数'', Amer. Journ. of Math. 4 (1881), p.~225-229.

{[}40{]} \_\_\_\_ ``论代数的相对形式'', Amer. Journ. of Math. 4 (1881), p.~221--225.

{[}41{]} H. POINCARÉ -- ``Sur les nombres complexes'', C. R. Acad. Sci. 99 (1884), p.~740--742; Œuvres, t. V, Paris (Gauthier-Villars), 1916--1954, p.~77--79.

{[}42{]} \_\_\_\_ ``Sur l'intégration algébrique des équations linéaires et les périodes des intégrales abéliennes'', Journ. de Math. (5) 9 (1903), p.~139--212; Œuvres, t. III, Paris (Gauthier-Villars), 1916--1954, p.~106--166.

{[}43{]} G. SCHEFFERS -- ``Zurückführung complexer Zahlensysteme auf typische Formen'', Math. Ann. 39 (1891), p.~293--390.

{[}44{]} I. SCHUR -- ``Über eine Klasse von Matrizen, die sich einer gegebenen Matrix zuordnen lassen'', Diss., Berlin, 1901; Gesammelte Abhandlungen, Band I, Berlin (Springer), 1973, p.~1--72.

{[}45{]} \_\_\_\_``Über die Darstellung der endlichen Gruppen durch gebrochene lineare Substitutionen'', Crelle's J. 127 (1904), p.~20--50; Gesammelte Abhandlungen, Band I, Berlin (Springer), 1973, p.~86--116.

{[}46{]} \_\_\_\_ ``Neue Begründung der Theorie der Gruppencharaktere'', Berliner Sitzungsber. (1905), p.~406--432; Gesammelte Abhandlungen, Band I, Berlin (Springer), 1973, p.~143--169.

{[}47{]} \_\_\_\_``Arithmetische Untersuchungen über endliche Gruppen linearer Substitutionen'', Berliner Sitzungsber. (1906), p.~164--184; Gesammelte Abhandlungen, Band I, Berlin (Springer), 1973, p.~177--197.

{[}48{]} J.-A. DE SÉGUIER - Éléments de la théorie des groupes abstraits, Paris (Gauthier-Villars), 1904.

{[}49{]} T. SKOLEM - Zur Theorie des assoziativen Zahlensysteme, Oslo Vid. Akad. Skr. I, 1927.

{[}50{]} J. M. WEDDERBURN - ``关于有限代数的一个定理'', Trans. Amer. Math. Soc. 6 (1905), p.~349-352.

{[}51{]} \_\_\_\_``论超复数'', Proc. Lond. Math. Soc. (2) 6 (1908), p.~77--118.

{[}52{]} \_\_\_\_``一类本原代数'', Trans. Amer. Math. Soc. 15 (1914), p.~162--166.

{[}53{]} \_\_\_\_ ``论可除代数'', Trans. Amer. Math. Soc. 22 (1921), p.~129--135.

{[}54{]} K. WEIERSTRASS -- ``Zur Theorie der aus n Haupteinheiten gebildeten complexen Grössen'', Gött. Nachr. (1884), p.~395--414; Mathematische Werke, t. II, Berlin (Mayer und Müller), 1895, p.~311--332.

\BourbakiBackMatter{记号索引}

\(A_{M}\) \ldots.. 1 K{[}G{]} \ldots.. 1 A{[}X{]} \ldots.. 9 A{[}X{]} \(_{\sigma,d}\) \ldots.. 11 {[}M:L{]} \ldots.. 32 IsA\{X,Y\} \(_{\ddagger}\) \ldots.. 51 cl(X) \ldots.. 51 F(A) \ldots.. 51 S(A) \ldots.. 51 s \(^{0}\) , d \(^{0}\) \ldots.. 53 T(V) \ldots.. 58 H(M) \ldots.. 58 αM \ldots.. 59 βV \ldots.. 59 MS \ldots.. 65 αM \ldots.. 69 τ(M) \ldots.. 80 S\_M \ldots.. 84 {[}B:A{]} \(_{s}\) \ldots.. 124 R\_A(M) \ldots.. 151 R(M) \ldots.. 151 R(A) \ldots.. 154 K(C) \ldots.. 186 {[}E{]}C,{[}E{]} \ldots.. 186 K'(C) \ldots.. 188 {[}E{]}'C \ldots.. 189 \(\mathrm{K}^{\prime}(\mathcal{C})^{+}\) 189 \(\ell_{\mathrm{S}}(\mathrm{E})\) 189 \(\mathcal{L}\mathcal{F}(\mathrm{A})\) 191 \(\mathrm{R(A)}\) 191 \(\mathrm{R_K(A)}\) 194 \(\mathrm{R_K(G)}\) 198 \(\mathcal{R}(\mathfrak{g})\) 199 \(\mathrm{K_0(A)}\) 199 \(i(f)\) 204 \(i(\mathrm{B,A})\) 204 \(h(f)\) 204 \(h(\mathrm{B,A})\) 204 \(\mathrm{C^n(A,P)}\) 239 \(\mathrm{H^n(A,P)}\) 240 \(\mathrm{A^f}\) 254 \(\mathrm{Iso}_{\mathrm{K}}\{\mathrm{A},\mathrm{B}\}\) 277 \(\operatorname {cl}(A)\) 277 \(\operatorname {cl}_{\mathrm{K}}(A)\) 277\\
Br(K) 278 \(\operatorname {Ex}_{\tau}(\mathrm{G},\mathrm{F})\) 286 \(\mathcal{I}_{\tau}\) 286 \(\mathrm{Z}^2 (\mathrm{G},\mathrm{F})\) 295 \(\mathrm{C}^1 (\mathrm{G},\mathrm{F})\) 295 \(\mathrm{H}^2 (\mathrm{G},\mathrm{F})\) 296 \(\operatorname {Res}_{\mathrm{G}'}^{\mathrm{G}}\) 299 \(\operatorname {Cor}_{\mathrm{H}}^{\mathrm{G}}\) 303

\(\mathrm{Coind}_{\mathrm{H}}^{\mathrm{G}}(\mathrm{E})\) 305 \(\mathbf{A}[\mathcal{E};\mathrm{L}]\) 315 \(\mathrm{H}^n (\mathrm{G},\mathrm{M})\) 324 \(\mathrm{Res}^q\) 325 \(\mathrm{Inf}^q\) 325 \(\mathrm{Cor}^q\) 326 \(\mathrm{Br(A)}\) 330 \(\mathrm{Br(B / A)}\) 330 \(\mathrm{Pcrd}_{\mathrm{A / K}}(a;\mathrm{X})\) 340 \(\mathrm{Trd}_{\mathrm{A / K}}\) 340 \(\mathrm{Nrd}_{\mathrm{A / K}}\) 340 \(\mathrm{Pm}_{\mathrm{A / K}}(a;\mathrm{X})\) 351 \(\Theta (\mathrm{A})\) 376 \(c_{\mathrm{E}}(x,x^{*})\) 377 \(\Theta^{\mathrm{ss}}(\mathrm{A})\) 380 \(\mathrm{Hom}_{\mathrm{G}}(\pi ,\pi ')\) 398

\(\mathcal{Z}_{\mathrm{K}}(\mathrm{G})\) 398 \(\chi_{\pi}\) 398 \(\pi \otimes \pi^{\prime}\) 399 \(\pi \boxtimes \pi^{\prime}\) 400 \(\operatorname{Res}_{\mathrm{H}}^{\mathrm{G}}(\pi)\) 402 \(\operatorname{Ind}_{\mathrm{H}}^{\mathrm{G}}(\sigma)\) 402 \(\operatorname{Coind}_{\mathrm{H}}^{\mathrm{G}}(\sigma)\) 403 \(\widehat{\mathbf{G}}\) 406 \(\mathrm{F}(\widehat{\mathrm{G}})\) 407 \(\det u\) 452 \(\det(A)\) 453 \(\mathbf{SL}_n(\mathrm{D})\) 455 \(\mathbf{SL}(\mathrm{V})\) 457 \(\operatorname{sdet} X\) 460 \(\operatorname{Hom}_{\mathrm{A}}^{\mathrm{f}}(\mathrm{E},\mathrm{F})\) 463 \(\operatorname{End}_{\mathrm{A}}^{\mathrm{f}}(\mathrm{E})\) 463

\BourbakiBackMatter{术语索引}

2-上边缘，295\\
2-闭上链，295

\BourbakiIndexLetter{A}

绝对半单代数，231 绝对半单模，229 绝对无根模，246 适配基，460 模的加性函数，183 模类的加性集，183 模的弱加性函数，183 伴随同构，59 拟逆元，441 代数 绝对半单代数，231 代数代数，359 阿廷代数，444 Azumaya 代数，271 中心代数，251 余诱导代数，305 伽罗瓦代数，310 (K, G)-代数，304 本原代数，443 四元数代数，361 半单代数，136 相似代数，280 单代数，120

阿廷\\
代数，444\\
模，1\\
环，4\\
自内射\\
环，53\\
Azumaya 代数，271\\
阿祖马亚定理，32

\BourbakiIndexLetter{B}

平衡模，77\\
基，适配，460\\
双交换子，77\\
双模\\
可逆，100\\
左，56\\
块，180\\
布饶尔群，279、330\\
伯恩赛德定理，84、429

\BourbakiIndexLetter{C}

分式演算，17\\
嘉当矩阵，201\\
凯莱（凯莱–哈密顿定理），338\\
中心\\
代数，251\\
函数，398\\
中心化子，77\\
特征标

布饶尔特征标，434\\
表示的特征标，382、398\\
谢瓦莱（谢瓦莱–瓦宁定理），355\\
模的类，51\\
上边缘，295\\
闭上链，295\\
与 \(\tau\)-扩张相伴的，296\\
系数\\
模的系数，377\\
表示的系数，377\\
余代数，余子代数，391\\
上同调\\
群，324\\
代数的上同调，239\\
余诱导\\
代数，305\\
表示，403\\
模的余长度，53\\
交换子，77\\
相容同态，324\\
分量，S 型同型分量，65\\
余限制，303、326\\
余转置，338\\
反模，8\\
覆盖，投射覆盖，161\\
判据，麦基判据，426\\
交叉乘积，315

\BourbakiIndexLetter{D}

可分解幂等元，146\\
分解，韦尔–菲廷分解，27\\
分式的定义，18\\
退化，非退化子代数，370\\
次数 左，124\\
表示的次数，373、397\\
子环上的次数，124、204\\
单代数的约化次数，253\\
稠密子环，129\\
稠密性定理，88\\
稠密性，雅各布森稠密性定理，442\\
导子，9\\
内导子，240\\
描述 典范描述，62、70

模的描述，69\\
同型模的描述，62\\
行列式，452、453\\
图，杨图，430\\
\(\tau\)-扩张的直像，289\\
可除，基数可除，123

\BourbakiIndexLetter{E}

元素 幂等元，145 整元，428 p-正则元，433 无重因子的元素，155 等价 代数的森田等价，100 模的森田等价，103 等价幂等元，40 本质子模，75 阶（指数），322 扩张 代数的扩张，248 扩张的乘积，293 分裂域扩张，281 τ-扩张，285 τ-扩张的乘积，293 半平凡 τ-扩张，286

\BourbakiIndexLetter{F}

因子，S 型单因子，141 代数的忠实表示，374 域 C1，358 有限域，357 分裂域，281 剩余域，26 菲廷（韦尔–菲廷分解），27 平坦，绝对平坦环，171 足，130 型，迹型，464 公式，弗罗贝尼乌斯公式，432 傅里叶反演公式，408 分式 分式演算，17 分式的定义，18 分式幺半群，17

分式环，18\\
弗罗贝尼乌斯\\
公式，432\\
互反律，203、403\\
环，53\\
函数\\
模的加性函数，183\\
中心函数，398\\
模的弱加性函数，183

\BourbakiIndexLetter{G}

伽罗瓦代数，310 自同构伽罗瓦群，274 伽罗瓦子环，274 盖尔范德（盖尔范德–马祖尔定理），368 生成模，79 戈尔迪定理，149 格罗滕迪克群，186、189 格罗滕迪克环，198 群 布饶尔群，279、330 自同构伽罗瓦群，274 格罗滕迪克群，186 直和的格罗滕迪克群，189 幺模群，455、457

\BourbakiIndexLetter{H}

哈密顿（凯莱–哈密顿定理），338\\
高度\\
同态的高度，204\\
子环的高度，204\\
模类的遗传集，183\\
希尔伯特\\
定理 90，327\\
零点定理，461\\
同态\\
相容同态，324\\
余限制同态，303、326\\
膨胀同态，299、325\\
连接同态，324\\
有限秩同态，463\\
限制同态，299、325\\
位似环，1

\BourbakiIndexLetter{I}

理想 极大理想，26、141、435 极小理想，50 幂零理想，149、156 素理想，133 本原理想，133 正则理想，435 半素双边理想，149 迹理想，80

幂等元 群代数的中心幂等元，408 可分解幂等元，146 不可分解幂等元，146

正交幂等元，146

不可分解幂等元，146 不可分解模，30

指数，322 同态的指数，204 子环的指数，204

诱导表示，402

膨胀同态，299、325

交结算子，398

逆像 闭上链的逆像，299 \(\tau\)-扩张的逆像，287

模的逆，100

可逆双模，100

同构，伴随同构，59

同型 S 型同型分量，65 同型模，61

\BourbakiIndexLetter{J}

雅各布森稠密性定理，88 雅各布森根，154 雅各布森定理，156

\BourbakiIndexLetter{K}

克鲁尔–雷马克–施密特定理，37

\BourbakiIndexLetter{L}

引理

中山引理，158\\
舒尔引理，47\\
长度\\
模的长度，72\\
环的长度，7\\
相连模，180\\
连接同态，324\\
具局部自同态环的模，31\\
局部环，26

\BourbakiIndexLetter{M}

麦基判据，426\\
马施克定理，401\\
矩阵，嘉当矩阵，201\\
极大\\
理想，26、141、435\\
子模，48\\
马祖尔（盖尔范德–马祖尔定理），368\\
平均，325\\
极小理想，50\\
模\\
绝对半单模，229\\
绝对无根模，246\\
阿廷模，1\\
平衡模，77\\
生成模，79\\
不可分解模，30\\
同型模，61\\
诺特模，1\\
表示的模，373\\
有限余长度的模，53\\
\(\mathcal{C}\) 型模，183\\
原始模，31\\
半阿廷模，74\\
半原始模，32\\
半单模，55\\
单模，45\\
具局部自同态环的模，31\\
模\\
相连模，180\\
稳定同构的模，200\\
分式幺半群，17\\
森田等价\\
代数的森田等价，100

模的森田等价，103\\
重数，72、189\\
原始重数，34\\
乘子\\
伪双模的乘子，445\\
代数的乘子，445

\BourbakiIndexLetter{N}

中山引理，158\\
诣零理想，156\\
幂零理想，156\\
幂零根，157\\
诺特（斯科伦–诺特定理），256\\
诺特模，1 诺特环，4\\
范数，约化范数，340\\
n-正则环，178

\BourbakiIndexLetter{O}

算子，交结算子，398 正交 正交幂等元，146 理想的正交，53 正交关系 特征标的正交关系，410 舒尔正交关系，410 特征标的第二正交关系，412

\BourbakiIndexLetter{P}

幂等元的分拆，146\\
置换表示，425\\
多项式\\
主多项式，351\\
约化特征多项式，340\\
素\\
素理想，133\\
素根，167\\
本原\\
本原理想，133\\
本原环，120\\
本原代数，443\\
原始模，31\\
主多项式，351\\
乘积

交叉乘积，315\\
外张量积，400\\
表示的张量积，399\\
投射覆盖，161\\
伪双模，445\\
伪模\\
左、右，439\\
半单伪模，442\\
单伪模，440\\
伪正则环，178

\BourbakiIndexLetter{Q}

拟单环，120\\
四元数，361

\BourbakiIndexLetter{R}

根 雅各布森根，154 模的根，151 环的根，154 代数的根，440 素根，167 互反律 弗罗贝尼乌斯互反律，203、403 约化 约化次数，253 正则理想，435 正则环，171 相关联的表示，416 关系 特征标的正交关系，410 舒尔正交关系，410 特征标的第二正交关系，412 表示 双正则表示，399 余诱导表示，403 完全可约表示，374 逆步表示，399 对偶表示，399 忠实表示，374 诱导表示，402 不可约表示，374 群的线性表示，397 代数的线性表示，373

矩阵表示，377\\
单项表示，427\\
置换表示，425\\
商表示，374\\
正则表示，375、399\\
半单表示，374\\
单表示\\
群的单表示，398\\
代数的单表示，374\\
转置表示，374\\
单位表示，399\\
酉表示，421\\
限制对偶，376\\
限制同态，299、325\\
环\\
绝对平坦环，171\\
阿廷环，4\\
自内射环，53\\
弗罗贝尼乌斯环，53\\
格罗滕迪克环，198\\
位似环，1\\
局部环，26\\
诺特环，4 \(n\)-正则环，178\\
分式环，18\\
多项式环，11\\
本原环，120\\
伪正则环，178\\
拟单环，120\\
半素环，149\\
半单环，136\\
单环，120\\
冯·诺伊曼正则环，171\\
无根环，154\\
佐恩环，171

\BourbakiIndexLetter{S}

饱和，左、右，19\\
舒尔正交关系，410\\
舒尔引理，47\\
特征标的第二正交关系，412\\
半阿廷模，74\\
半素环，149\\
半原始模，32

半单代数，136 半单模，55 半单伪模，442 半单环，136

半单化，191 半平凡 τ-扩张，286 模类的集 加性，183 遗传，183 相似代数，280 单代数，120 S 型单因子，141 单伪模，440 群的单表示，398 单环，120 单模，45 斯科伦（斯科伦–诺特定理），256 底座，65、130 右底座、左底座，130 分裂域，281 子代数 极大交换子代数，261 极大交换半单子代数，264 极大平展子代数，264 非退化子代数，370 余子代数，391 子模 本质子模，75 极大子模，48 子伪模，439 子表示，374 子环 稠密子环，129 伽罗瓦子环，274 弱伽罗瓦子环，274 半单模的支集，66、190

\BourbakiIndexLetter{T}

定理

阿祖马亚定理，32\\
伯恩赛德定理，84、429\\
凯莱–哈密顿定理，338\\
谢瓦莱–瓦宁定理，355\\
盖尔范德–马祖尔定理，368\\
戈尔迪定理，149\\
希尔伯特定理 90，327\\
希尔伯特零点定理，461\\
雅各布森定理，156\\
雅各布森稠密性定理，88、442\\
克鲁尔–雷马克–施密特定理，37\\
马施克定理，401\\
斯科伦–诺特定理，256\\
曾定理，358、360\\
韦德伯恩定理，120、135、245、357\\
迹\\
迹型，464\\
迹理想，80\\
表示的迹，398\\
约化迹，340\\
迹型映射，115\\
平延，457\\
曾定理，358、360

\BourbakiIndexLetter{U}

幺模群，455\\
单位元，左、右，435

\BourbakiIndexLetter{W}

瓦宁（谢瓦莱–瓦宁定理），355 弱伽罗瓦子环，274 韦德伯恩定理，120、135、245、357 韦尔（韦尔–菲廷分解），27

\BourbakiIndexLetter{Y}

杨图，430

\BourbakiIndexLetter{Z}

佐恩环，171
